\section{Wyniki}
\label{sec:wyniki-testow}

Testy uruchomiłam 26 sierpnia 2026 na Node.js 24.13.0 pod Windows 11, w Vitest 2.1.9 i Playwright 1.62.1.

\begin{table}[htbp]
\centering
\caption{Wyniki uruchomienia testów -- opracowanie własne}
\label{tab:wyniki-uruchomienia}
\begin{tabular}{|p{0.23\textwidth}|p{0.20\textwidth}|c|p{0.19\textwidth}|c|}
\hline
\textbf{Poziom} & \textbf{Polecenie} & \textbf{Przyp.} & \textbf{Wynik} & \textbf{Czas} \\
\hline
Jednostkowy + API & \texttt{npm test} w \texttt{backend} & 107 & 107 zaliczonych & 3,0 s \\
\hline
Jednostkowy + komponenty & \texttt{npm test} w \texttt{frontend} & 107 & 107 zaliczonych & 13,2 s \\
\hline
Row Level Security & \texttt{npm run test:rls} & 16 & pominięte & -- \\
\hline
Systemowy & \texttt{npm run test:e2e} & 7 & pominięte & -- \\
\hline
\textbf{Razem} & & \textbf{237} & \textbf{214 zaliczonych, 23 pominięte} & \textbf{16,2 s} \\
\hline
\end{tabular}
\end{table}

Rozkład testów na pliki pokazuje tabela~\ref{tab:rozklad-plikow}. Widać w nim kształt piramidy z podrozdziału~\ref{sec:strategia-testow}: na dole 107 testów jednostkowych, na górze siedem ścieżek systemowych.

\begin{table}[htbp]
\centering
\caption{Rozkład testów na pliki -- opracowanie własne}
\label{tab:rozklad-plikow}
\begin{tabular}{|p{0.18\textwidth}|p{0.54\textwidth}|c|}
\hline
\textbf{Poziom} & \textbf{Plik} & \textbf{Testy} \\
\hline
Jednostkowy & \texttt{frontend/tests/unit/monument-price.test.ts} & 29 \\
Jednostkowy & \texttt{frontend/tests/unit/safe-path.test.ts} & 26 \\
Jednostkowy & \texttt{frontend/tests/unit/photo-crop.test.ts} & 17 \\
Jednostkowy & \texttt{backend/tests/unit/http-helpers.test.js} & 17 \\
Jednostkowy & \texttt{frontend/tests/unit/password-policy.test.ts} & 15 \\
Jednostkowy & \texttt{frontend/tests/unit/translations.test.ts} & 14 \\
Jednostkowy & \texttt{frontend/tests/unit/i18n-context.test.tsx} & 11 \\
Jednostkowy & \texttt{backend/tests/unit/exchange-rate.test.js} & 10 \\
Jednostkowy & \texttt{frontend/tests/unit/api-client.test.ts} & 9 \\
Jednostkowy & \texttt{frontend/tests/unit/currency.test.tsx} & 7 \\
\hline
API & \texttt{backend/tests/integration/auth.routes.test.js} & 25 \\
API & \texttt{backend/tests/integration/contact-and-public.test.js} & 25 \\
API & \texttt{backend/tests/integration/order-submit.test.js} & 17 \\
API & \texttt{backend/tests/integration/admin.test.js} & 14 \\
API & \texttt{backend/tests/integration/require-auth.test.js} & 14 \\
API & \texttt{backend/tests/integration/rate-limit.test.js} & 9 \\
\hline
Baza & \texttt{backend/tests/rls/policies.test.js} & 16 \\
\hline
Komponenty & \texttt{frontend/tests/components/designer-page.test.tsx} & 9 \\
Komponenty & \texttt{frontend/tests/components/protected-route.test.tsx} & 7 \\
Komponenty & \texttt{frontend/tests/components/catalog-page.test.tsx} & 6 \\
Komponenty & \texttt{frontend/tests/components/contact-form.test.tsx} & 6 \\
\hline
Systemowy & \texttt{frontend/tests/e2e/system.spec.ts} & 7 \\
\hline
\end{tabular}
\end{table}

\subsection{Testy pominięte}
\label{ssec:przypadki-niewykonane}

23 testy zostały pominięte, bo oba poziomy korzystające z prawdziwej bazy potrzebują więcej niż zwykłego komputera do pracy. Testy RLS zakładają i kasują prawdziwe konta oraz rekordy, więc włączają się dopiero po ustawieniu zmiennej \texttt{RLS\_TEST\_ENABLED} i wskazaniu testowego projektu Supabase. Ścieżki systemowe potrzebują uruchomionej aplikacji, serwera i bazy oraz kont montera i administratora założonych wcześniej.

Wynik należy więc czytać tak: 214 testów wykonanych i zaliczonych oraz 23 gotowe do uruchomienia w środowisku testowym. Do tego jeden test RLS jest oznaczony jako taki, który ma się nie udać, bo dokumentuje lukę opisaną w podrozdziale~\ref{sec:testy-rls}. W pełnym środowisku wyjdzie z niego 15 zaliczonych i jeden nieudany, a nie 16 zaliczonych.

\subsection{Błędy znalezione przez testy}
\label{ssec:usterki}

Testy przydały się jeszcze zanim uruchomiłam cały zestaw, bo pokazały dwa błędy, których nie wychwyciłabym czytaniem kodu.

Pierwszy dotyczył samych testów komponentów. Na komputerze bez pliku konfiguracyjnego aplikacji klienckiej przechodziły, ale na moim, gdzie ten plik wskazuje adres uruchomionego serwera, 14 z 32 testów nie przechodziło -- żądania omijały atrapy i leciały pod prawdziwy adres. Wynik testu zależał więc od ustawień komputera, wbrew zasadzie z podrozdziału~\ref{sec:strategia-testow}. Naprawiłam to, wymuszając w konfiguracji testów pusty adres backendu.

Drugi wyszedł przy przypadku z tabeli~\ref{tab:przypadki-wyceny}, czyli materiale bez podanej ceny. Zakładałam, że funkcja pominie ten składnik i pokaże cenę niepełną, a okazało się, że zwraca wartość nieliczbową, która trafiłaby do interfejsu w miejsce kwoty. Błąd nie ujawniłby się przy normalnym korzystaniu, bo wszystkie materiały w katalogu mają cenę, ale wyszedłby przy pierwszym niekompletnym rekordzie. Dodałam sprawdzenie, czy cena jest liczbą.
